%Syftet med studien är att undersöka möjligheterna att utveckla kommunikation och informationsutbyte mellan lokförare med hjälp av C-DAS, bortom både det grundläggande protokoll som Trafikverket tillämpar och som i sin tur bygger på vad som fastställs enligt SFERA\@.
\paragraph{}
%\hspace{-11pt}ptptSyftet med studien är att utvärdera och vidareutveckla den av Trafikverket fastställda digitala kommunikationen och informationsutbytet mellan lokförare och tågklarerare samt undersöka eventuella förbättringsmöjligheter avseende C-DAS-kommunikationens innehåll som ännu ej prövats.
%\hspace{-11pt}ptptSyftet med studien är att utvärdera fastställd digital kommunikation och informationsutbyte mellan lokförare och tågklarerare inom C-DAS samt undersöka eventuella förbättringsmöjligheter avseende C-DAS-kommunikationens innehåll som ännu ej prövats om det kan identifieras.
\hspace{-11pt}Syftet med studien är att undersöka i vilken utsträckning fastställd digital kommunikations- och informationsutbyte mellan lokförare och tågklarerare tillgodoser lokförares informations- och kommunikationsbehov.
\paragraph{}
\hspace{-11pt}För att den nya paradigmen med C-DAS/TMS ska lyckas behöver en så stor del som möjligt av tågtrafiken framföras med aktiv C-DAS\@.
Det ligger därför i Trafikverkets intresse att järnvägsföretagen tillämpar en C-DAS-lösning i sin verksamhet.
%Om datautbytet mellan lokförare och tågklarerare kan förbättras så att förarnas förtroende för C-DAS ökar skulle det kunna leda till en högre grad av efterlevnad av hastighetsanvisningar hos förarna.
%Om datautbytet mellan lokförare och tågklarerare kan förbättras så att förarnas förtroende för C-DAS ökar skulle det kunna leda till en högre grad av efterlevnad av hastighetsanvisningar hos förarna.
Om datautbytet mellan lokförare och tågklarerare tillämpas så att förarnas informations- och kommunikationsbehov inte tillgodoses skulle det kunna leda till lägre förtroende och en lägre grad av efterlevnad av hastighetsanvisningar hos förarna.
Omvänt skulle ett tillgodosett informationsbehov med högt förtroende hos förarna i sin tur kunna leda till högre grad av efterlevnad av hastighetsanvisningar.
I förlängningen skulle detta kunna leda till snabbare införandetakt av C-DAS hos järnvägsföretagen samt good-will för Trafikverket.
\begin{enumerate}
    \item Hur ser benägenheten ut för lokförare att följa hastighetsrekommendationer från en DAS-applikation idag?
    \item Hur väl tillgodoser befintlig standard för datautbyte inom C-DAS (SFERA) lokförares informationsbehov avseende efterlevnad av hastighetsanvisningar?
%   \item Hur ser informationsbehovet ut för lokförare gällande följandet av C-DAS-anvisningar jämfört med det som erbjuds enligt SFERA-standarden?
%   \item Tillgodoses detta behov av SFERA-standarden idag?
    \item Hur kan kommunikation- och informationsutbyte mellan lokförare och tågklarerare förbättras genom att gå utöver det grundläggande protokoll Trafikverket tillämpar idag (SFERA)?
%   \item Vilka data från lokförare till TMS som inte täcks in i dagens informationskedja kan stärka tågklarerarnas beslutsstöd och i förlängningen operativa förmåga?
    \item Hur kan gamification användas för att stödja och motivera lokförare i deras användande av C-DAS-verktyg?
\end{enumerate}